\documentclass[a4paper]{article}
% Español (México) con XeLaTeX: fontspec para fuentes Unicode, polyglossia
% para la división silábica y los nombres del documento en la variante mexicana.
\usepackage{fontspec}
\usepackage{polyglossia}
\setdefaultlanguage[variant=mexican]{spanish}
% Los nombres chinos van como «pinyin (original)»: xeCJK da a los caracteres
% Han su propia fuente; sin ella XeLaTeX los omite con un aviso.
% FandolSong no cubre algunos caracteres de nombres (U+73FA); WenQuanYi Zen Hei sí.
\usepackage[AutoFallBack=true]{xeCJK}
\setCJKmainfont{FandolSong-Regular.otf}
\setCJKfallbackfamilyfont{\CJKrmdefault}{WenQuanYi Zen Hei}
% polyglossia no traduce los nombres de listings.
\AtBeginDocument{\providecommand{\lstlistingname}{}\renewcommand{\lstlistingname}{Listado}%
  \providecommand{\lstlistlistingname}{}\renewcommand{\lstlistlistingname}{Índice de listados}}
\usepackage{amsmath, amssymb}
\usepackage{graphicx}
\usepackage[margin=2.35cm]{geometry}
\usepackage[most]{tcolorbox}
\usepackage{booktabs}
\usepackage{tabularx}
\usepackage{longtable}
\usepackage{array}
\usepackage{float}
\usepackage{xcolor}
\usepackage{hyperref}
\usepackage{fancyhdr}

\hypersetup{colorlinks=true, linkcolor=blue!55!black, urlcolor=blue!60!black}
\graphicspath{{./}{frames/}{frames/candidates/}{figures/}}
\setlength{\parskip}{0.55em}
\setlength{\parindent}{2em}
\setlength{\emergencystretch}{3em}
\renewcommand{\arraystretch}{1.18}

\newtcolorbox{knowledgebox}[1]{enhanced, colback=blue!5!white, colframe=blue!70!black, colbacktitle=blue!70!black, coltitle=white, fonttitle=\bfseries, title={#1}, sharp corners, breakable}
\newtcolorbox{importantbox}[1]{enhanced, colback=yellow!10!white, colframe=yellow!75!black, colbacktitle=yellow!75!black, coltitle=black, fonttitle=\bfseries, title={#1}, sharp corners, breakable}
\newtcolorbox{warningbox}[1]{enhanced, colback=red!5!white, colframe=red!70!black, colbacktitle=red!70!black, coltitle=white, fonttitle=\bfseries, title={#1}, sharp corners, breakable}

\pagestyle{fancy}
\fancyhf{}
\fancyfoot[C]{\thepage}
\fancyfoot[R]{\footnotesize\color{gray}\href{https://github.com/hqhq1025/ai-course-notes}{AI Course Notes}}
\renewcommand{\headrulewidth}{0pt}
\renewcommand{\footrulewidth}{0pt}

\newcommand{\videourl}{https://www.youtube.com/watch?v=Xxz5uh0L1mE}
\newcommand{\repourl}{https://github.com/hqhq1025/ai-course-notes}

\begin{document}
\begin{titlepage}
\centering
{\Large Notas de entrevista de tipo panorama técnico\par}
\vspace{1.0cm}
{\huge\bfseries Historia técnica de los Agent, el OpenClaw Moment y la disolución de las fronteras\par}
\vspace{0.6cm}
{\Large Zhang Xiaojun conversa con Su Yu (苏煜)\par}
\vspace{0.4cm}
{\large Elaborado a partir del video público de Zhang Xiaojun Podcast, la transcripción hecha con GPU y fotogramas clave locales\par}
\vspace{0.3cm}
{\large 2026-05-01\par}
\vspace{0.9cm}
\includegraphics[width=0.88\textwidth,height=0.42\textheight,keepaspectratio]{cover.jpg}\par
\vfill
\begin{tcolorbox}[width=0.92\textwidth, colback=black!2!white, colframe=black!55, sharp corners]
\textbf{Título del video}: 139. [Panorama de los Agent] Con Su Yu sobre la historia técnica de los Agent, el OpenClaw Moment, la disolución de las fronteras y su irradiación en la sociedad\par
\textbf{Canal del video}: Zhang Xiaojun Podcast\par
\textbf{Duración del video}: 02:17:49\par
\textbf{Enlace del video}: \href{\videourl}{\nolinkurl{\videourl}}\par
\textbf{Descripción de los materiales}: YouTube no tiene subtítulos disponibles; el texto principal se elaboró a partir de la transcripción con faster-whisper large-v3 en CUDA, los capítulos del video, la portada, los fotogramas clave y figuras didácticas propias.\par
\textbf{Página del proyecto}: \href{\repourl}{\nolinkurl{\repourl}}\par
\end{tcolorbox}
\end{titlepage}

\tableofcontents
\newpage
\section{Introducción: por qué este episodio es una clase sobre la genealogía técnica de los Agent}

Este episodio no es una entrevista común con un invitado, sino una historia técnica de los Agent contada de viva voz. El invitado, Su Yu (苏煜), es profesor del Departamento de Ciencias de la Computación de la Universidad Estatal de Ohio y fundador de NeoCognition; desde hace años investiga temas como Language Agent, Computer Use Agent, Mind2Web, LM Planner y MMMU. La entrevista intenta responder una gran pregunta: por qué en 2026 el relato sobre la IA pasa del Chat al Agent, por qué un producto como OpenClaw provocó en la industria una sacudida parecida al ChatGPT Moment, y por qué se están reorganizando las fronteras entre browser, desktop, mobile, GUI, CLI, API y coding.

\begin{importantbox}{La pregunta central de este episodio}
Si el ChatGPT Moment se entiende como la confirmación pública del paradigma de los LLM, el OpenClaw Moment se parece más a la confirmación pública del paradigma de los Agent: el modelo ya no solo responde preguntas, sino que entra en un entorno, usa herramientas, ejecuta acciones, aprende de forma continua y se acerca poco a poco a un universal digital agent.
\end{importantbox}




\subsection{Resumen de la sección}

Este episodio se presta para organizarse como una nota sobre la «genealogía técnica de los Agent»: primero se define qué es un Agent; después se repasan el logical agent, el neural agent, el semantic parsing y el language agent; por último se discuten OpenClaw, el universal digital agent, el continual learning, el world model y la difusión en la industria.
\section{La definición mínima de Agent: entidad, entorno y actividad orientada a metas}

La definición de Agent que da Su Yu (苏煜) es muy sencilla: un Agent es una entidad con límites definidos que opera en un entorno externo y realiza goal-directed activities, es decir, actividades orientadas a metas. La definición es amplia a propósito, porque animales, personas, robots, agentes que operan páginas web y coding agents caben en este marco. Lo decisivo no es si se parece a un humano, sino si tiene límites identificables, entradas del entorno y un comportamiento impulsado por metas.

\begin{figure}[H]
\centering
\includegraphics[width=0.92\textwidth]{agent-loop.png}
\caption{El ciclo mínimo de un Agent: entorno, observación, razonamiento, acción y memoria.}
\end{figure}

\begin{knowledgebox}{Lectura de la figura: cómo leer el ciclo del Agent}
En el extremo izquierdo de la figura está el entorno, que puede ser una página web, un escritorio, un repositorio de código, una base de datos o el mundo de un robot; arriba, Observation es el estado percibido; en el centro, el Agent razona a partir de la meta, la observación y la memoria; a la derecha, Action modifica el entorno; abajo, Memory/World Model consolida la experiencia de interacción en un estado reutilizable. La conclusión que sustenta es esta: un Agent no es una respuesta única, sino un sistema de lazo cerrado.
\end{knowledgebox}

\begin{importantbox}{Una fórmula simplificada}
Una acción de un Agent puede escribirse como:
$$
a_t = \pi(o_t, m_t, g)
$$
donde $a_t$ es la acción del paso $t$; $o_t$ es la observación actual; $m_t$ es la memoria o el modelo del mundo; $g$ es la meta; $\pi$ es la política o función de decisión. Esta fórmula no aparece explícitamente en la entrevista; es una condensación didáctica de la definición de Su Yu.
\end{importantbox}

\subsection{Por qué la definición debe ser amplia}

Si la definición es demasiado estrecha, es fácil confundir un Agent con una forma de producto concreta, como un agente web, un agente de escritorio o un coding agent. La definición de Su Yu considera todos estos casos como instancias de una misma clase de sistema en entornos distintos. Browser use, desktop use, mobile use, GUI, CLI, API y coding son solo means to an end; la meta real es un universal digital agent capaz de completar todo tipo de tareas en el digital world.

\begin{warningbox}{Malentendido frecuente: un Agent no es una ventana de chat con plugins}
Que una ventana de chat invoque una herramienta es apenas un corte temprano de lo que es un Agent. Un Agent real debe percibir el estado, planificar acciones, ejecutarlas, aprovechar la retroalimentación y mantener la meta a lo largo de tareas de varios pasos. Reducir un Agent a «un LLM que sabe invocar APIs» subestima la importancia de memory, world model, reliability y cost.
\end{warningbox}

\subsection{Resumen de la sección}

La definición mínima de Agent es: una entidad con límites definidos que actúa en un entorno con base en metas. Esta definición permite colocar en una misma línea histórica los primeros sistemas simbólicos, el aprendizaje por refuerzo, el semantic parsing, el uso de herramientas por LLM y los coding agents modernos.
\section{Historia de la evolución técnica: de Logical Agent a Language Agent}

Su Yu (苏煜) divide la historia técnica de los Agent en varias etapas: los primeros logical agent, los neural agent posteriores al año 2000, la rama de semantic parsing, que va del lenguaje a la semántica formal, y los language agent de los últimos tres años. Esta genealogía importa porque evita ver el auge de los Agent en 2026 como algo nuevo surgido de la nada.

\begin{figure}[H]
\centering
\includegraphics[width=0.96\textwidth]{agent-history-timeline.png}
\caption{Genealogía técnica de los Agent: de la lógica simbólica al OpenClaw Moment.}
\end{figure}

\begin{knowledgebox}{Lectura de la figura: la historia técnica no es una historia lineal de victorias}
La línea de tiempo comienza con los logical agent, pasa por los neural agent y el semantic parsing, y termina en los language agent. La tendencia central no es que «los métodos nuevos sustituyen por completo a los antiguos», sino que los LLM convierten el lenguaje en una nueva interfaz general, de modo que los problemas que antes estaban dispersos entre la lógica, la planeación, la invocación de herramientas y la interacción con el entorno vuelven a converger.
\end{knowledgebox}

\subsection{Logical Agent: lógica simbólica y sistemas expertos}

Los Logical Agent eran la visión dominante de la IA temprana: describir el mundo con lógica formal, reglas y sistemas expertos, y luego tomar decisiones mediante un sistema de inferencia. Su ventaja es que son interpretables y tienen una estructura clara; su debilidad es que el mundo real es demasiado complejo, y la cobertura de reglas, el manejo de excepciones y la entrada perceptual son difíciles de escalar. Su Yu enfatiza que perseguir un Agent de inteligencia artificial completo resultaba algo excesivo con las condiciones técnicas de la época, y que eso terminó por dividir la IA en subcampos como la visión, el procesamiento de lenguaje natural y el razonamiento lógico.

\noindent\begin{tabularx}{\textwidth}{p{0.22\textwidth}p{0.34\textwidth}X}
\toprule
Etapa & Mecanismo & Limitaciones \\
\midrule
Logical Agent & Lenguajes lógicos, sistemas expertos, planeación simbólica & Interpretables pero frágiles; dependen de reglas escritas a mano y difícilmente cubren entornos abiertos. \\
Neural Agent & Redes neuronales, aprendizaje por refuerzo, modelos de percepción & Mayor capacidad perceptual, pero carecen de una interfaz de lenguaje general y de capacidad para usar herramientas complejas. \\
Semantic Parsing & Del lenguaje natural a SQL, expresiones lógicas, llamadas a API & Eficaz en entornos específicos, pero cada entorno debe modelarse por separado. \\
Language Agent & LLM como modelo del mundo lingüístico e interfaz de herramientas & Gran capacidad de generalización, pero reliability, memory, costo y aprendizaje continuo siguen siendo cuellos de botella. \\
\bottomrule
\end{tabularx}

\subsection{Semantic Parsing: la prehistoria de los Language Agent}

El Semantic Parsing intenta mapear el lenguaje natural a acciones formalizadas, por ejemplo consultas SQL, consultas a grafos de conocimiento o llamadas a API. Se parece mucho a los Language Agent modernos: ambos deben convertir el lenguaje en conductas ejecutables. La diferencia es que, antes de los LLM, los sistemas por lo general solo funcionaban en una base de datos, un sitio web o un grafo de conocimiento específicos; con los LLM, el modelo incorpora conocimientos previos lingüísticos y conocimiento del mundo más sólidos, lo que le permite generar conductas razonables en muchos más entornos.

\begin{knowledgebox}{Términos clave: Semantic Parsing y Language Agent}
El Semantic Parsing resuelve el problema de «convertir una oración en una forma ejecutable», por ejemplo, transformar «buscar la población de cierta ciudad» en SQL. Los Language Agent extienden esta idea a entornos abiertos: el modelo no solo genera semántica formal, sino que también planea, invoca herramientas, observa resultados y sigue actuando. No hay una ruptura entre ambos, sino una relación de prehistoria y extensión.
\end{knowledgebox}

\subsection{Language Agent: por qué el lenguaje es un acelerador}

Su Yu usa la evolución humana como analogía del papel del lenguaje: el lenguaje apareció muy tarde en la historia evolutiva del ser humano, pero aceleró enormemente el desarrollo de la civilización. De manera similar, los LLM dan a los Agentes de IA una interfaz simbólica general. El lenguaje natural, los lenguajes de programación, los diagramas y los gestos pueden considerarse formas amplias de language; un programming language es solo un tipo de language más formalizado y más adecuado para que una máquina lo ejecute.

\begin{importantbox}{Tesis central: el lenguaje no es una interacción superficial, sino una interfaz con el modelo del mundo}
Con los LLM, los Agent ya no necesitan construir desde cero un analizador semántico para cada entorno. El lenguaje se convierte en la capa intermedia que conecta objetivos, planes, herramientas, retroalimentación y memoria. Coding es importante precisamente porque es uno de los lenguajes formales más potentes del digital world.
\end{importantbox}

\subsection{Resumen de la sección}

La historia técnica de los Agent no comenzó en 2026. Lo nuevo es que los LLM convierten el language en una interfaz general, de modo que los objetivos de los logical agent, la ejecutabilidad del semantic parsing, la capacidad perceptual de los neural agent y la invocación de herramientas vuelven a converger.
\section{Los Language Agents de los últimos tres años: ReAct, planificación, uso de herramientas y la ola del código abierto}

Según la entrevista, en los últimos tres años los Language Agents avanzaron más rápido que en las décadas anteriores. Entre los hitos importantes están ReAct, LLM Planner, Mind2Web, Toolformer y AutoGPT, además de una serie de trabajos sobre web agents, computer use agents y robot planning. La transcripción reconoce Toolformer como «2-former»; por el contexto, se refiere a la línea de uso de herramientas del tipo Toolformer de Meta. Esta nota lo trata según el sentido de la entrevista.

\begin{figure}[H]
\centering
\includegraphics[width=0.92\textwidth]{language-agent-stack.png}
\caption{Pila de un Language Agent: del objetivo del usuario a las herramientas, la retroalimentación del entorno y la memoria.}
\end{figure}

\begin{knowledgebox}{Lectura de la figura: los niveles de la pila de un Language Agent}
En el nivel superior está el objetivo del usuario; debajo, el lenguaje natural, los lenguajes de programación y las instrucciones estructuradas; en medio, la planificación y el razonamiento del LLM; más abajo, las interfaces externas: herramientas, API, GUI, CLI, bases de código, etc.; en la base, la retroalimentación del entorno y la memoria. La conclusión que sustenta es esta: un Language Agent no es un truco de prompt, sino una pila de ingeniería de ciclo cerrado.
\end{knowledgebox}

\subsection{ReAct: un insight sencillo pero de gran alcance}

La esencia de ReAct es alternar reasoning y acting: el modelo primero razona, luego actúa, luego observa y vuelve a razonar. Su Yu (苏煜) subraya que muchos trabajos sobre Agents parecen técnicamente sencillos, pero proponer la abstracción correcta en el momento correcto es muy difícil. El valor de ReAct está en que llevó al LLM de la generación de texto en un solo paso a un ciclo de interacción con el entorno.

\noindent\begin{tabularx}{\textwidth}{p{0.22\textwidth}p{0.34\textwidth}X}
\toprule
Trabajo/línea & Problema que resuelve & Relación con el hilo de este episodio \\
\midrule
ReAct & Hacer que el modelo alterne reasoning y acting & Establece el patrón básico de interacción entre el LLM y el entorno. \\
LLM Planner & Usar el LLM para planificación robótica o embodied planning & Aplica el modelo de lenguaje a planes de acción, no solo a conversar. \\
Mind2Web & Construir un benchmark de web/computer use agents & Convierte las tareas web en una capacidad de Agent evaluable. \\
Toolformer & Hacer que el modelo aprenda cuándo invocar herramientas & Conecta el modelo de lenguaje con el ecosistema de API/herramientas externas. \\
AutoGPT & Proyecto temprano de long-horizon agent de código abierto & Muestra lo que el público imagina de los agentes autónomos y también expone sus problemas de fiabilidad. \\
\bottomrule
\end{tabularx}

\subsection{De la proof of concept a la etapa intensiva en recursos}

Su Yu repasa cómo pasó de semantic parsing a language agents y explica por qué cada vez más investigadores dejan la universidad para emprender. Los primeros trabajos sobre Agents eran más bien proofs of concept: una buena idea podía demostrarse con poco costo. A partir de 2025, las ideas de Agents verdaderamente interesantes requieren cada vez más GPU, API, equipos de ingeniería y capacidad de probar y corregir con rapidez, algo que no encaja del todo con la estructura de recursos de la universidad.

\begin{warningbox}{No confundir «el proyecto de código abierto se volvió popular» con «el problema está resuelto»}
Proyectos como AutoGPT y OpenClaw atraen la atención con rapidez porque muestran lo que es posible. Pero lo posible no equivale a un producto fiable. La gestión del estado, la recuperación de errores, el control de costos y los límites de seguridad en tareas de largo plazo son la parte verdaderamente difícil de los Agents.
\end{warningbox}

\subsection{Resumen de la sección}

El desarrollo de los Language Agents en los últimos tres años es el paso de la generación de texto en un solo paso al ciclo cerrado con el entorno. ReAct, la planificación, el uso de herramientas, los web agents y los agentes autónomos de código abierto impulsaron en conjunto la llegada del OpenClaw Moment.
\section{El OpenClaw Moment: por qué se parece al ChatGPT Moment}

El conductor compara el auge de OpenClaw con ChatGPT, y Su Yu (苏煜) también ve semejanzas entre ambos: el ChatGPT Moment marcó el reconocimiento público del paradigma de los LLM, mientras que el OpenClaw Moment marca el reconocimiento público de un paradigma de Agent más automatizado y personalizado. En ninguno de los dos casos la tecnología de base surgió de repente ese día; fueron capacidades que ya existían y que se mostraron en la forma de producto adecuada.




\subsection{Qué significa «moment»}

Un moment no es la primera vez que un proyecto logra cierta capacidad, sino el momento en que se forma de golpe un consenso social. Antes de ChatGPT ya había modelos de lenguaje, y antes de OpenClaw ya había web agent, tool use, planning y coding agent. Pero cuando llega el moment, el capital, los productos, la investigación y las expectativas de los usuarios se reorganizan con rapidez.

\begin{figure}[H]
\centering
\includegraphics[width=0.94\textwidth]{agent-productization-funnel.png}
\caption{Del Moment al sistema en producción: el embudo de conversión de Agent en producto.}
\end{figure}

\begin{knowledgebox}{Lectura de la figura: qué capas quedan por atravesar después del OpenClaw Moment}
La figura se estrecha capa por capa: Moment, Demo, Benchmark, Deployment y Operating System. Muestra que el éxito repentino de un producto es solo el punto de partida del consenso; la verdadera ventaja competitiva proviene de capacidades de sistema que puedan evaluarse, desplegarse, aprender a largo plazo y dar lugar a un ecosistema.
\end{knowledgebox}


\begin{importantbox}{La esencia del OpenClaw Moment}
El OpenClaw Moment no es «un proyecto de código abierto que se volvió muy popular», sino el momento en que la industria empieza a creer que un modelo puede trabajar en tareas de horizonte largo a través de herramientas, interfaces y tareas distintas. Convierte al universal digital agent de un problema de investigación en una idea de producto.
\end{importantbox}

\subsection{China y Estados Unidos difieren en su modo de difusión}

En la entrevista se menciona que el pattern de difusión tecnológica difiere entre China y Estados Unidos. En China, la capa de aplicaciones se mueve más rápido y llega a toda la población; en Estados Unidos, la difusión suele empezar en las empresas, entre los desarrolladores y en el software de productividad. Esto significa que el alcance social de los Agent no depende solo de la capacidad de los modelos, sino también del ecosistema de productos, la composición de los usuarios, las compras corporativas, la cultura de los desarrolladores y la velocidad de creación de empresas.

\begin{knowledgebox}{Dos vías de difusión industrial}
La vía estadounidense suele ser enterprise/productivity first: primero entra en el desarrollo, la oficina y los procesos empresariales, y después se extiende gradualmente hacia fuera. La vía china puede ser más consumer/application first: la escala de usuarios, las plataformas de contenido, la propagación en redes sociales y el relato de la superaplicación pesan más. Por ello, la forma de producto de los Agent puede ser distinta en cada lado.
\end{knowledgebox}

\subsection{Resumen de la sección}

La importancia del OpenClaw Moment está en el giro del consenso: los Agent dejan de ser solo una demo de investigación y se convierten en la puerta de entrada a la siguiente generación de flujos de trabajo digitales. Se parece al ChatGPT Moment no porque la tecnología sea la misma, sino porque cambió las expectativas de la industria.
\section{La disolución de las fronteras: Browser, Desktop, Mobile, GUI, CLI, API y Coding}

Su Yu (苏煜) insiste en que, en sus inicios, el campo de los Agent distinguía entre browser use, desktop use, mobile use, GUI, text-based representation, CLI, API, coding, etc., pero que estas divisiones son provisionales. Lo que finalmente se busca es un universal digital agent. Todas estas interfaces son medios, no el fin último.

\begin{figure}[H]
\centering
\includegraphics[width=0.92\textwidth]{boundary-convergence.png}
\caption{La disolución de las fronteras: múltiples puntos de entrada a entornos digitales están convergiendo hacia un universal digital agent.}
\end{figure}

\begin{knowledgebox}{Lectura de la figura: por qué el coding es el núcleo de la disolución de las fronteras}
En la figura, todos los puntos de entrada se conectan con el universal digital agent. Browser, desktop, mobile y GUI son la interacción de superficie; CLI, API y coding son interfaces de control más formalizadas. Su Yu considera que el coding es el fabric del digital world, porque muchas interfaces pueden, en última instancia, expresarse, renderizarse o manipularse mediante código.
\end{knowledgebox}

\subsection{La GUI no desaparecerá, y la CLI tampoco lo dominará todo}

La entrevista no llega al extremo de afirmar que «la CLI reemplazará por completo a la GUI». Hay dos razones. Primera: en el mundo real existe una gran cantidad de legacy system, por ejemplo en bancos, software empresarial y sistemas antiguos, que no se reescribirán con rapidez. Segunda: aunque la CLI sea el óptimo global para los Agent, eso no significa que sea la mejor opción en todos los escenarios locales. Muchas GUI existentes ya son good enough para las personas y las organizaciones.

\begin{warningbox}{El óptimo local es distinto del óptimo global}
Una interfaz técnicamente más adecuada para los Agent no necesariamente sustituirá de inmediato a las interfaces existentes. Los costos de migración de las empresas, los hábitos de los usuarios, la regulación, los sistemas heredados y las cuentas económicas determinan la trayectoria real. El juicio técnico sobre los Agent debe considerarse junto con el entorno de despliegue.
\end{warningbox}

\subsection{Un Programming Language también es un Language}

El conductor propone la metáfora de que «el lenguaje natural es el andamiaje de los humanos y el coding es el andamiaje de las máquinas». Su Yu añade que language nunca se ha limitado al lenguaje natural: los lenguajes de programación, los diagramas y los gestos son language en sentido amplio. Un programming language es un formal language, más preciso y más ejecutable, y por ello resulta especialmente importante cuando un Agent opera en el digital world.

\subsection{Resumen de la sección}

El desenlace de los Agent no es la victoria individual del browser agent, el desktop agent o el coding agent, sino la convergencia de múltiples puntos de entrada en el nivel de la tarea. La importancia del coding radica en que unifica muchas interfaces como objetos que pueden expresarse, combinarse y ejecutarse.
\section{Continual Learning, World Model y Expert Agent}

Cuando el conductor preguntó cuál es el mayor cuello de botella de los Agent, Su Yu (苏煜) presentó memory, self learning, continual learning, world model, specialization y expert agent como distintas facetas de una misma cuestión. Lo que hoy les falta a los Agent es aprender de forma continua a partir de la interacción y convertir la experiencia en capacidades estables.

\begin{figure}[H]
\centering
\includegraphics[width=0.94\textwidth]{continual-learning-map.png}
\caption{Relación entre Continual Learning, World Model y Expert Agent.}
\end{figure}

\begin{knowledgebox}{Lectura de la figura: por qué estos términos forman una cadena}
La experiencia de la izquierda proviene de la interacción con el entorno; el continual/self learning actualiza esa experiencia y la convierte en capacidades; el world model representa el aprendizaje de las regularidades del entorno; la specialization representa la formación de un experto de dominio; el resultado final es un Agent más confiable, más rápido y de menor costo. La conclusión que sustenta es que memory, world model y la especialización no son roadmaps independientes entre sí, sino una misma cadena de acumulación de capacidades.
\end{knowledgebox}

\subsection{Términos clave: vocabulario de los cuellos de botella de los Agent}

\noindent\begin{tabularx}{\textwidth}{p{0.20\textwidth}p{0.32\textwidth}X}
\toprule
Término & Problema que resuelve & Lugar en este episodio \\
\midrule
Memory & Conservar el historial y las preferencias relevantes en tareas de varios pasos & Sin memoria, el Agent actúa cada vez como novato y tiende a repetir errores. \\
Continual Learning & Actualizar las capacidades de forma continua a partir de nuevas interacciones & Hace que el Agent no solo ejecute guiones, sino que acumule experiencia. \\
World Model & Aprender el estado del entorno, la causalidad y qué acciones son posibles & Sustenta la planificación, la predicción de consecuencias y la operación segura. \\
Specialization & Convertirse en expert agent de un dominio & Pasar de asistente general a trabajador vertical confiable. \\
Reliability & Tasa de finalización estable de las tareas & Uno de los mayores problemas actuales en la conversión en producto de los Agent. \\
Cost Effectiveness & Si el costo por tarea es aceptable & Determina si el Agent puede pasar de la demo a producción. \\
\bottomrule
\end{tabularx}

\begin{importantbox}{Formulación unificada del mayor cuello de botella}
El mayor cuello de botella de los Agent no es un módulo de memory aislado, un prompt aislado ni una herramienta aislada, sino la incapacidad de transformar, como lo hace una persona, la experiencia de largo plazo en capacidades profesionales estables. Continual learning, world model y specialization sirven, en última instancia, a la confiabilidad, la velocidad y el costo.
\end{importantbox}

\subsection{La señal de los forward-deployed engineers}

En la entrevista se menciona que algunas empresas recurren a forward-deployed engineers: envían ingenieros a las instalaciones del cliente para ayudarle a construir sus Agents. Esto indica que los Agent todavía no se han difundido con una barrera de entrada baja: hace falta entender los procesos del cliente, el entorno de herramientas, los modos de falla y las restricciones de seguridad. El producto sigue en una etapa con un alto componente de servicio.

\begin{warningbox}{La dificultad del despliegue no se debe solo a un modelo débil}
Que a los Agent les cueste entrar en las empresas no se debe solo a la insuficiencia de capacidades del modelo, sino también a la comprensión de los procesos, la gestión de permisos, la integración de datos, la responsabilidad por los errores, los límites de seguridad y el cálculo de costos. Atribuir todas las fallas a que «el modelo no es lo bastante inteligente» lleva a juzgar mal la dificultad de la conversión en producto.
\end{warningbox}

\subsection{Resumen de la sección}

La línea principal de la siguiente etapa de los Agent es el aprendizaje continuo y el modelado del mundo. El progreso real se manifestará en confiabilidad, rapidez, bajo costo y especialización, y no solo en completar secuencias de acciones más vistosas en una demo.
\section{Las bets de las grandes empresas tecnológicas y el emprendimiento: por qué los investigadores fundan empresas}

Su Yu (苏煜) considera que antes las bets de cada empresa en torno a los Agents eran muy distintas entre sí, pero que ahora tienden a coincidir, porque el camino de Anthropic se convirtió en el modelo a seguir para la industria. Tanto Anthropic como OpenAI están convergiendo hacia la dirección de la productivity; Google cuenta con modelos sólidos y con una posición privilegiada en su ecosistema, aunque la repercusión de sus productos no necesariamente está a la altura; xAI, entre otras, también presta atención al computer user agent o al general digital agent.




\subsection{La estructura de recursos de la universidad, las empresas y el emprendimiento}

Su Yu explica por qué pasó de la universidad al emprendimiento: la universidad es adecuada para las weird ideas, la proof of concept y los marcos conceptuales; pero cuando los Agents entran en una etapa intensiva en recursos, en la que se necesitan GPU, API, un equipo sólido y una ejecución de ingeniería rápida, la estructura de recursos de la universidad ya no se ajusta. Emprender no es simplemente una cuestión de dinero, sino una forma de seguir investigando; lo que cambió es la forma que adopta la investigación.

\noindent\begin{tabularx}{\textwidth}{p{0.18\textwidth}p{0.36\textwidth}X}
\toprule
Ámbito & Para qué es adecuado & Para qué no es adecuado \\
\midrule
Universidad & Marcos conceptuales, validación de bajo costo, weird ideas de largo plazo & Ingeniería a gran escala, iteración rápida de productos, experimentos que exigen muchos recursos. \\
Grandes empresas tecnológicas & Entrenamiento de modelos grandes, infraestructura, integración del ecosistema & La exploración libre en múltiples direcciones puede verse limitada por los objetivos de la organización. \\
Startup & Ensayo y error rápido, enfoque en el producto, construcción de sistemas completos en torno a los Agents & Fuerte presión de recursos; es necesario encontrar un wedge claro y una vía comercial. \\
\bottomrule
\end{tabularx}

\subsection{El valor del conceptual framework}

Al final de la entrevista, Su Yu dice que le gusta construir conceptual frameworks: no es que tenga una memoria especialmente buena ni que reaccione con especial rapidez, sino que puede aprender muchas cosas, enlazarlas y ver las conexiones entre ellas. Eso es también lo más valioso de esta revisión de conjunto: coloca la historia de los Agents, el lenguaje, las herramientas, el coding, el continual learning, el world model y la difusión en la industria dentro de un marco unificado.

\begin{importantbox}{Por qué vale la pena convertir este episodio en una nota de alto nivel}
Muchos podcasts son densos en opiniones pero de estructura laxa; este episodio es justo lo contrario, porque en sí mismo construye un marco conceptual. Lo importante al organizarlo no es repetir cada ejemplo, sino conservar este marco: la definición de Agent, la historia técnica, la pila del Language Agent, la disolución de fronteras, el cuello de botella del aprendizaje continuo y las bets de la industria.
\end{importantbox}

\subsection{Resumen de la sección}

Los Agents pasaron de la proof of concept a una etapa intensiva en recursos. La división del trabajo entre la universidad, las grandes empresas tecnológicas y las startups está cambiando; quienes sean capaces de construir marcos conceptuales y, a la vez, movilizar recursos de ingeniería tendrán más facilidad para impulsar la siguiente etapa de los Agents.
\section{Evaluación, confiabilidad y seguridad: de la demo al sistema en producción}

Esta entrevista no deja el problema del Agent en la pregunta «¿se puede hacer una demo impresionante?». Su Yu (苏煜) reconduce varias veces la discusión hacia reliability, speed, cost effectiveness y deployment. Una demo de agent puede mostrar que «sabe hacerlo» mediante un entorno, unos ejemplos y unas indicaciones humanas diseñados con cuidado; pero un sistema en producción tiene que responder otras preguntas: qué tasa de fallos tiene ante tareas desconocidas, si los fallos son recuperables, si el costo es controlable, si los límites de permisos están claros y cómo se acotan las fugas de datos o las operaciones erróneas.

\begin{importantbox}{Cuatro umbrales para pasar de la demo a producción}
El primero es la tasa de éxito de la tarea, no solo que cada acción individual sea correcta; el segundo es la recuperabilidad, es decir, si tras un fallo se puede localizar el error y continuar; el tercero es el costo, si los token, las llamadas a herramientas y el respaldo humano de las tareas de largo alcance resultan económicos; el cuarto es la seguridad: cuando el agent puede operar sobre entornos reales, los permisos, la auditoría y el peor escenario posible tienen que incorporarse desde el diseño.
\end{importantbox}

\subsection{Por qué el benchmark no es suficiente}

Los primeros benchmark de web agent o de computer-use permitieron a los investigadores comparar sistemas, pero con frecuencia no cubren por completo un despliegue real. En el entorno real de una empresa hay sesiones iniciadas, permisos, datos históricos, procesos organizacionales, sistemas heredados, requisitos de auditoría e instrucciones incompletas de los usuarios. Que un agent tenga un buen desempeño en un benchmark estándar no significa que pueda operar de forma segura en un banco, en el área legal, en el sector salud, en una cadena de herramientas de investigación y desarrollo o en los sistemas internos de una empresa.

\noindent\begin{tabularx}{\textwidth}{p{0.20\textwidth}p{0.34\textwidth}X}
\toprule
Dimensión de evaluación & Qué hay que medir & Por qué importa \\
\midrule
Task success & Si la tarea de largo alcance se completa al final & Evita optimizar solo un clic aislado o una respuesta única. \\
Recovery & Si tras un error puede autodiagnosticarse, revertir y reintentar & Determina si el sistema puede operar sin supervisión. \\
Cost & Token, llamadas a herramientas, tiempo de espera, intervención humana & Determina si el agent puede desplegarse a escala. \\
Safety & Límites de permisos, acciones peligrosas, privacidad y auditoría & Condición necesaria para pasar de la demo a la operación real del negocio. \\
Adaptation & Si puede aprender de entornos nuevos y de retroalimentación nueva & Vincula el continual learning con el expert agent. \\
\bottomrule
\end{tabularx}

\subsection{La diferencia entre Security y Safety}

En la entrevista aparece una distinción que se pasa por alto con facilidad: muchos problemas de safety son, en el fondo, problemas de capacidad, porque el agent, como un principiante, no entiende el entorno y comete errores elementales; en cambio, la security se orienta más al worst-case scenario y requiere métodos específicos. Es decir, aumentar la capacidad puede reducir muchos errores por descuido, pero no sustituye el control de permisos, los entornos aislados, la auditoría ni los ejercicios de equipo rojo.

\begin{warningbox}{Mejorar la capacidad no es diseñar la seguridad}
Un agent más inteligente quizá haga menos clics equivocados, pero también puede ser más capaz de eludir las restricciones. Un sistema en producción debe diseñar por separado la mejora de capacidad y los mecanismos de seguridad: la capacidad se encarga de completar la tarea; la seguridad, de acotar el espacio de la tarea, registrar el comportamiento y bloquear las acciones inaceptables.
\end{warningbox}


\subsection{Lista práctica: qué preguntar antes de construir un sistema de Agent}

Si el contenido de esta entrevista se convierte en una checklist de ingeniería, hay que plantear al menos ocho preguntas. Primera: cuáles son los límites del Agent, quién puede asignarle objetivos y quién puede detenerlo. Segunda: cuál es el entorno al que se enfrenta y cómo se observa el estado de ese entorno. Tercera: qué acciones puede realizar y cuáles deben prohibirse o requerir confirmación humana. Cuarta: cómo se determina automáticamente el éxito de la tarea y qué tareas requieren aceptación humana. Quinta: cómo se recupera tras un fallo y si puede revertir cambios o generar registros de auditoría. Sexta: cómo acumula experiencia y si esa experiencia puede contaminar tareas posteriores. Séptima: si su costo es predecible. Octava: desde qué punto de entrada del producto lo invocarán los usuarios de forma natural.

\begin{importantbox}{Criterio de ingeniería}
Que un sistema de Agent pueda ponerse en producción no depende de que complete una tarea asombrosa en una demostración, sino de que las ocho preguntas anteriores puedan responderse de forma estable. Esta lista también sirve como marco reutilizable para leer otras entrevistas sobre Agent.
\end{importantbox}

\subsection{Resumen de la sección}

Llevar un Agent a producción no es solo un problema del modelo. La evaluación, la recuperación, el costo, la seguridad y la adaptación al entorno determinan en conjunto si puede pasar de la demo a un sistema real. Por eso el continual learning y el world model se vuelven el eje central: en última instancia deben servir a la confiabilidad, no solo a demostraciones más vistosas.
\section{Irradiación social: cómo el Agent pasa de la investigación a la aplicación masiva}

Otra palabra clave del OpenClaw Moment es «irradiación social». Su Yu (苏煜) y el conductor coinciden en que China y Estados Unidos difunden la tecnología a ritmos distintos: en Estados Unidos es más fácil entrar primero en productivity, enterprise y herramientas para desarrolladores; en China, en cambio, la capa de aplicaciones podría masificarse más rápido. Esta diferencia implica que una misma tecnología de Agent dará lugar a formas de producto distintas en mercados distintos.

\subsection{Por qué coding es el primer frente principal}

Coding es un buen candidato para ser el primer frente principal de los Agent por razones directas: las tareas son verificables, el entorno está relativamente formalizado, la retroalimentación se obtiene de la compilación, las pruebas, el lint y los resultados de ejecución, y los usuarios están dispuestos a pagar por una mayor productividad. En comparación, un agent de consumo de propósito general tiene que resolver muchos más problemas de necesidades ambiguas, preferencias personales, privacidad y diseño de interacción.

\begin{knowledgebox}{El desbordamiento de coding hacia el universal digital agent}
Si el coding agent funciona, su efecto no se limita a los programadores. El código es la expresión de base de muchos sistemas digitales; un agent capaz de escribir código, modificar configuraciones, llamar API, generar scripts y leer registros se extenderá de forma natural hacia el análisis de datos, la operación, la automatización de oficina, las herramientas internas y los procesos empresariales.
\end{knowledgebox}

\subsection{Diferencias en los patrones de difusión entre China y Estados Unidos}

El mercado de software empresarial en Estados Unidos es grande y la vía de pago directo es clara, por lo que los Agent tienden a implementarse primero en coding, ventas, atención al cliente, oficina y enterprise workflow. En China, el ecosistema de aplicaciones para el consumidor final es complejo, y las plataformas, el contenido, las redes sociales, el comercio electrónico y los sistemas de recomendación están más integrados entre sí; por eso podrían surgir más rápido puntos de entrada de agent masivos, ligados a escenarios concretos, orientados al entretenimiento o con forma de superaplicación.

\begin{importantbox}{La forma del producto depende de la estructura social}
Una misma capacidad técnica no genera por sí sola el mismo producto. En Estados Unidos, el Agent podría presentarse primero como herramienta de productividad empresarial; en China, podría entrar más rápido en el contenido, las redes sociales, el comercio electrónico y los asistentes personales. La ruta tecnológica debe entenderse junto con la estructura del mercado.
\end{importantbox}

\subsection{Resumen de la sección}

La irradiación social de los Agent no es una simple difusión tecnológica. La transforman los modelos de negocio, los hábitos de los usuarios, las compras empresariales, los ecosistemas de plataformas y las formas de difusión cultural. Solo al entender esto se puede explicar por qué una misma ola tecnológica avanza a ritmos distintos en China y en Estados Unidos.
\section{Conceptual Framework: cómo condensar el contenido de este episodio en un mapa mental}

Al final de la entrevista, Su Yu (苏煜) dijo que le gusta build conceptual framework. Esa expresión también sirve para entender el contenido de este episodio: un Agent no es un producto aislado, sino un conjunto de conexiones entre conceptos. Parte del ideal de «agente» de la IA temprana, pasa por la lógica simbólica, las redes neuronales y el análisis semántico, y con la llegada de los LLM recupera una interfaz unificada; después, apoyado en coding, el uso de herramientas, world model y continual learning, empieza a entrar en los sistemas de producción.

\begin{importantbox}{Un mapa mental expresado en palabras}
Este episodio puede condensarse en cuatro capas: la primera es la definición: un Agent es una entidad con límites que actúa dentro de un entorno orientada a objetivos; la segunda es la historia: Logical Agent, Neural Agent, Semantic Parsing y Language Agent son etapas distintas de un mismo problema; la tercera es la interfaz: language/coding/tool/API/GUI/CLI están convergiendo; la cuarta es la puesta en producción: la confiabilidad, la velocidad, el costo, la seguridad y el aprendizaje continuo determinan si puede implantarse en la práctica.
\end{importantbox}

\subsection{Por qué la «disolución de fronteras» no es un eslogan}

La disolución de fronteras no significa que todas las interfaces vayan a desaparecer, sino que una misma tarea puede expresarse mediante interfaces distintas. Las operaciones en el navegador pueden expresarse con el DOM y con scripts, las operaciones de escritorio pueden expresarse con UI automation, las llamadas a API pueden encapsularse en código, y el coding, a su vez, puede generar y reescribir esas interfaces. Lo que un Agent necesita aprender de verdad es la semántica de la tarea, el estado del entorno y el espacio de acciones posibles, en lugar de quedar atado a una forma de entrada concreta.

\subsection{Del problema de investigación al problema de producto}

En la etapa de investigación suele preguntarse «¿puede este agent hacerlo?»; en la etapa de producto hay que preguntarse «¿puede hacerlo de forma estable, barata y segura?». Por eso este episodio vuelve una y otra vez a reliability, speed, cost y forward deployment. El OpenClaw Moment permitió que todos vieran lo que es posible, pero un sistema de producción necesita que esa posibilidad atraviese la organización, los permisos, los datos, los sistemas heredados y los hábitos de los usuarios.

\begin{warningbox}{Mal uso del marco conceptual}
Un marco conceptual no sirve para todo. Meter todo en un solo marco tiende a borrar las diferencias. El valor del marco de este episodio está en explicar tendencias, pero cada producto concreto debe volver a la tarea, los datos, los usuarios, el entorno y el costo. En distintos escenarios, la importancia relativa de GUI, CLI, API y coding será diferente.
\end{warningbox}

\subsection{Resumen de la sección}

Lo más sólido de este episodio no es una predicción aislada, sino un marco: el Agent es la nueva unificación de un problema histórico, language es la interfaz, coding es la capa formal del mundo digital, continual learning es la vía de acumulación de capacidades de la siguiente etapa, y la puesta en producción depende de la confiabilidad y el costo.
\section{Lecciones para la cola de trabajo siguiente: por qué EP139 debe servir de modelo para los episodios de revisión}

Este episodio es de tipo review/survey, no una entrevista a una persona. Ofrece a la cola de trabajo siguiente de Zhang Xiaojun AI una plantilla de procesamiento: cuando hay slides, se hace slide-complete; cuando no hay slides, no se insertan una y otra vez fotogramas de las personas, sino que se generan diagramas conceptuales, glosarios y diagramas de mecanismos. En los programas de revisión, el valor de la nota está en reconstruir la estructura del conocimiento, no en condensar la conversación.

\begin{knowledgebox}{Plantilla de generación para pódcasts de revisión}
Primero, extraer la genealogía técnica; segundo, identificar los términos centrales; tercero, convertir los juicios expresados de forma oral en tablas o diagramas de flujo; cuarto, señalar las controversias y los límites; quinto, usar la sección de síntesis para enlazar con el siguiente episodio o el siguiente tema. Solo así la nota resultante se parece a un material de clase y no a una transcripción editada.
\end{knowledgebox}

\subsection{Resumen de la sección}

La metodología de EP139 se aplicará a la cola de trabajo siguiente sobre IA e internet en sentido amplio: en el cuerpo del texto solo entran la portada y las imágenes que de verdad aportan información; cuando no hay contenido visual, la estructura didáctica se apoya en gráficos generados; en cada episodio se conservan transcript, manifest, coverage y visual QA.


\section{Apéndice: índice de términos y ruta de repaso}

\subsection{Índice de términos}

\begin{longtable}{p{0.23\textwidth}p{0.32\textwidth}p{0.36\textwidth}}
\toprule
Término & Explicación en una frase & Función en este episodio \\
\midrule
Agent & Entidad con límites que actúa en un entorno orientada a objetivos & Objeto central de todo el texto; abarca personas, animales, agentes de software y agentes del mundo digital. \\
Logical Agent & Agente inteligente basado en reglas lógicas y razonamiento simbólico & Muestra que Agent no es un concepto nuevo, sino un objetivo temprano de la IA. \\
Semantic Parsing & Convertir lenguaje natural en representaciones ejecutables como SQL, API o expresiones lógicas & Antecedente directo de Language Agent. \\
Language Agent & Agent que usa un LLM como interfaz de lenguaje y puede planificar y usar herramientas & Base técnica del OpenClaw Moment. \\
ReAct & Patrón que alterna Reasoning y Acting & Lleva al LLM de la respuesta única a un ciclo de interacción con el entorno. \\
Tool Use & El modelo invoca herramientas, API o funciones externas & Permite al modelo superar el límite de la respuesta de solo texto. \\
Computer Use Agent & Agent capaz de operar navegadores, escritorios, dispositivos móviles o entornos de código & Puerta de entrada importante hacia el universal digital agent. \\
World Model & Modelo interno del estado del entorno, de su causalidad y de las acciones posibles & Base del continual learning y de la confiabilidad. \\
Continual Learning & Actualizar las capacidades de forma continua a partir de nuevas interacciones & Candidato a eje principal del desarrollo de los Agent en 2026. \\
Expert Agent & Agent especializado que acumula habilidades en un dominio & Forma resultante que resuelve los problemas de confiabilidad, velocidad y costo. \\
\bottomrule
\end{longtable}

\subsection{Ruta de repaso}

Si solo se quiere repasar rápidamente este episodio, se puede leer en cuatro pasos: primero, la sección 2, para dominar la definición mínima de Agent; después, las secciones 3--4, para entender la historia técnica y la pila de Language Agent; luego, las secciones 5--7, para entender OpenClaw, la disolución de fronteras y el continual learning; por último, las secciones 8--10, para entender las bets de las grandes empresas tecnológicas, la evaluación en producción y la difusión social. Así, la pregunta «¿por qué los Agent se volvieron importantes de repente?» deja de ser una noticia de productos y se replantea como un problema de sistemas técnicos.

\begin{knowledgebox}{Conexión entre este episodio y el siguiente}
EP139 presenta el marco técnico de los Agent; la entrevista con Luo Fuli (罗福莉) en EP138 sitúa ese marco dentro de los problemas de un laboratorio de modelos: el postentrenamiento, la RL infra, la asignación de GPU y la igualdad dentro de la organización. Leer primero EP139 y después EP138 facilita entender por qué «el paradigma de los Agent depende mucho del postentrenamiento».
\end{knowledgebox}
\section{Síntesis y ampliación}

\subsection{Conclusiones centrales de este episodio}

\begin{enumerate}
\item Agent no es un concepto nuevo: la IA imagina Agents desde sus primeras etapas. Lo nuevo es que los LLM hacen del lenguaje una interfaz general para usar herramientas.
\item Language Agent hereda de semantic parsing la capacidad de producir algo ejecutable, pero gracias a los LLM generaliza mucho mejor.
\item OpenClaw Moment se parece a ChatGPT Moment: en ambos casos, una forma de producto amplificó capacidades que ya existían y eso generó un consenso social.
\item Las fronteras entre browser, desktop, mobile, GUI, CLI, API y coding se están disolviendo; coding es la capa de expresión básica del digital world.
\item Los cuellos de botella centrales de Agent pueden unificarse como continual learning y world model: aprender de la experiencia y formar capacidades expertas confiables.
\item En 2026, la competencia entre Agents no es solo entre modelos: también se compite en producto, despliegue, costo, confiabilidad y ejecución organizacional.
\end{enumerate}

\subsection{Preguntas abiertas}

\begin{itemize}
\item ¿La ruta dominante de continual learning será el aprendizaje basado en world model o sistemas más ingenieriles de memoria, recuperación y flujos de trabajo?
\item ¿Universal digital agent madurará primero en coding/productivity o tendrá su auge en las aplicaciones de consumo?
\item ¿Las fronteras entre GUI, CLI, API y coding se unificarán en la capa del modelo o se mantendrán por mucho tiempo varios puntos de entrada al producto?
\item ¿El modelo de forward-deployed engineer será una forma transitoria o la norma a largo plazo para implantar Agents en las empresas?
\end{itemize}

\subsection{Lecturas adicionales}

\begin{itemize}
\item \textit{Artificial Intelligence: A Modern Approach}: punto de entrada clásico para entender la IA temprana y la definición de Agent.
\item Trabajos como ReAct, Toolformer, AutoGPT, Mind2Web y LM Planner: para entender los hitos clave de Language Agent en los últimos tres años.
\item Los benchmarks de Computer Use Agent, Web Agent, GUI Agent y Coding Agent: para entender por qué es difícil evaluar un universal digital agent.
\end{itemize}

\begin{importantbox}{Síntesis final}
El objetivo último de Agent no es «conversar mejor», sino «acumular mejor experiencia en el entorno y cumplir objetivos». La importancia de OpenClaw Moment está en que permitió a la industria ver esa posibilidad; lo difícil ahora es convertirla en sistemas confiables, de bajo costo y desplegables.
\end{importantbox}

\end{document}
